\documentclass[a4paper, 12pt]{article}

\usepackage[english, russian]{babel}
\usepackage[T2A]{fontenc}
\usepackage[utf8]{inputenc}
\setlength{\parindent}{1cm}
\usepackage{graphicx}
\usepackage{amsmath}
\usepackage{amssymb}
\usepackage{amsfonts}
\usepackage{array}
\usepackage{booktabs}
\usepackage{wrapfig}
\usepackage{caption} %заголовки плавающих объектов
\captionsetup[table]{justification=centering} %установки для заголовков таблиц
\captionsetup[figure]{justification=centering} %установки для заголовков таблиц



\begin{document}
	\begin{center}
		\subsubsection*{Расчётное задание по определению оценок вероятностных характеристик случайных величин с использованием конечной выборки наблюдений}
		\subsubsection*{Выполнил студент второго курса Физико-механического интститута, высшей школы МПУ: Куликовских Илья Михайлович}
		\subsubsection*{группа: 5031503/30001}
		\subsubsection*{Преподаватель: Ермакова Наталья Юрьевна}
	\end{center}
	
	\quad \textbf{Задание №1.} Дана выборка объема n = 100. Требуется проверить гипотезу о равномерном
	распределении генеральной совокупности.
	Для этого:
	\begin{enumerate}
		\item Найти выборочные оценки математического ожидания и среднеквадратического отклонения.
		\item Построить гистограмму относительных частот, подобрав количество и ширину каналов гистограммы, обеспечивающих корректное представление случайной величины. Представить таблицу для группированного статистического ряда выборки.
		\item На том же графике представить полигон приведенных частот.
		\item Подобрать параметр функции плотности вероятности. Построить график функции плотности вероятности.
		\item Используя критерий согласия Пирсона, проверить гипотезу о принадлежности выборки равномерно распределенной генеральной совокупности на уровне значимости $\alpha$= 0,05. Привести таблицу для критерия $\chi^2$.
		\item Привести график экспериментальной и теоретической функций распределения.
	\end{enumerate}
	
	\quad \textbf{Решение.} Для начала определим значения оценок математического ожидания и дисперсии для данной выборки:
	
	\begin{center}
		$ \hat{\mu}= \dfrac{1}{n} \sum\limits_{i=1}^n {x_i} = 0.038286 $
		
		$ x_{max} = 7.8057, \text{"  "} x_{min} = -7.7984 $
		
		$ \hat{\sigma^2} = \dfrac{1}{n-1} \sum\limits_{i=1}^{n} (x_i - \left< x \right>)^2 = 19.9084051117212 $
		
		\begin{table}[h!]
		\centering
		\begin{tabular}{|c|c|c|c|c|c|c|}
			\hline
				$x$ & $(\langle x \rangle -x_i)$ & $(\langle x \rangle-x_i)^2$ \\
			\hline
			-0.8992 & -0.937486 & 0.878880000196 \\
			-0.3132 & -0.351486 & 0.123542408196 \\
			6.1919 & 6.153614 & 37.866965260996 \\
			-7.7984 & -7.836686 & 61.413647462596 \\
			1.2559 & 1.217614 & 1.482583852996 \\
			-6.937 & -6.975286 & 48.654614781796 \\
				$\vdots$ & $\vdots$ & $\vdots$ \\
				-0.2302 & -0.268486 & 0.072084732196 \\
				-1.8436 & -1.881886 & 3.541494916996 \\
				-7.2377 & -7.275986 & 52.939972272196 \\
			\hline
				$\frac{1}{n} \sum\limits_{i=1}^n {x_i} = 0.038286$ & $\sum\limits_{i=1}^n (x_i-\left< x \right>) = 0$ & $\frac{1}{n-1} \sum\limits_{i=1}^{n} (x_i - \left< x \right>)^2 = 19.9084051117212$ \\
			\hline
		\end{tabular}
		\caption{определение математического ожидания, линейного отклонения, дисперсии}
	\end{table}
	
	\end{center}
	
	\quad После опредления основных параметров генеральной совокупности необходимо построить гисторгамму относительных частот. Для этого определим оптимальную ширину канала и посчитаем частоту попадания в канал. Согласно лекционному материалу оптимальное количество каналов примерно равно $\sqrt{n}$, тогда для данной выборки количество каналов гистограммы будет равно 10. Ширина канала определяется как: 
	
	\begin{center}
		$ h = \dfrac{x_{\text{max}}-x_{\text{min}}}{10} = \dfrac{7.8057 - (-7.7984)}{10} = \dfrac{15.6041}{10} = 1.56041 $
	\end{center}
	
	\begin{table}[h!]
		\centering
		\begin{tabular}{|c|c|c|c|c|}
			\hline
			№ & [$x_{k-1}$, $x_k$) & $n_i$ & $\frac{n_i}{n}$ & $\frac{n_i}{nh}$ \\
			\hline
			1 & [-7.7984, -6.2380) & 9 & 0.09 & 0.05768 \\
			2 & [-6.2380, -4.6776) &11 & 0.11 & 0.0705 \\
			3 & [-4.6776, -3.1172) & 10 & 0.1 & 0.0641 \\
			4 & [-3.1172, -1.5568) & 7 & 0.07 & 0.0449 \\
			5 & [-1.5568, 0.0036) & 17 & 0.17 & 0.1089 \\
			6 & [0.0036, 1.5640) &6 & 0.06 & 0.0385 \\
			7 & [1.5640, 3.1244) & 9 & 0.09 & 0.0577 \\
			8 & [3.1244, 4.6848) & 10 & 0.1 & 0.0641 \\
			9 & [4.6848, 6.2452) & 11 & 0.11 & 0.0705 \\
			10 & [6.2452, 7.8057] & 10 & 0.1 & 0.0641 \\
			\hline
		\end{tabular}
		\caption{Построение гистограммы относительных частот}
	\end{table}

	\quad Теперь построим функцию плотности вероятности для равномерного распределения. В качестве параметра возьмём $x_{max} - x_{min}$. Данный выбор обуславливается следующим образом:
	
	\begin{center}
	
		$f(x) =
			\begin{cases} 
				0, \forall x \in (-\infty, x_{min}) \cap (x_{max}, +\infty)\\
				 \dfrac{1}{x_{\text{max}} - x_{\text{min}}}, \forall x \in [x_{min}, x_{max}]
			\end{cases} \implies$
			
			$\implies f(x) = 
			\begin{cases}
				0, \forall x \in (-\infty, x_{min}) \cap (x_{max}, +\infty) \\
				\dfrac{1}{1.56041}, \forall x \in [x_{min}, x_{max}]
			\end{cases}$
			
	
		Проверим условие нормировки:
		
		$ \displaystyle\int\limits_{-7.7984}^{7.8057} {\dfrac{1}{15.6041}} dx = \dfrac{x}{15.6041} \bigg|_{-7.7984}^{7.8057} = \dfrac{7.8057 - (-7.7984)}{15.6041} = \dfrac{15.6041}{15.6041} = 1$
	\end{center}
	
	\quad Гистограмма и полигон относительных частот, а также функция плотности вероятности изображены на рисунке 1. Далее проверим гипотезу о равномерном распределении случайной величины. Для этого воспользуемся критерием $\chi^2_\alpha$ на уровне значимости $\alpha = 0.05$.
	
	\begin{figure}[h!]
		\centering
		\includegraphics[width=1\textwidth]{../images/гистограмма.png}
		\caption{ Гистограмма относительных частот, полигон относительных частот, функция плотности вероятности.}
	\end{figure}
	
	\quad Теоретические частоты для равномерного распределения определяются как: 
	
	\begin{center}
		$n_i ' = \dfrac{n}{10} = 10 \implies \chi^2 = \sum\limits_{i=0}^{n} \dfrac{(n_i - n_i ')^2}{n_i '} = 8.1078$
	
	\end{center}
	
	\quad Для нахождения критического значения $\chi^2_{\alpha}$ необходимо определить число степеней свободы. В данном случае оно  определяется количеством каналов гистограммы относительных часот за вычетом количества параметров в распределении:
	
	\begin{center}
	
		$ k = \sqrt{n} - 2 = 10 -2 = 8 \implies $
	
		$\implies$Критическое значение: $\chi^2_\alpha = 15.5 \implies 8.1078 < 15.5 \implies $ Гипотеза принимается.
		
		Функция распределения вероятности имеет вид: 
		
		$ F(x) = 
		\begin{cases}
			0, \forall x \in (- \infty, x_{min}) \\
			\dfrac{x}{1.56041}, \forall x \in [x_{min}, x_{max}] \\
			1, \forall x \in (x_{max}, +\infty) 
		\end{cases}$
		
		\begin{figure}[h!]
			\centering
			\includegraphics[width=1\textwidth]{../images/распределение.png}
			\caption{ Теоритическая и эксперементальная функции распределения вероятности.}
		\end{figure}
		
	\end{center}
	
	
	\quad \textbf{Задание №2.} Дана выборка объема n = 100. Требуется проверить гипотезу о нормальном
	распределении генеральной совокупности. Для этого:
	\begin{enumerate}
		\item Найти выборочные оценки математического ожидания и среднеквадратического
		отклонения.
		\item Построить гистограмму относительных частот, подобрав количество и ширину
		каналов гистограммы, обеспечивающих корректное представление случайной
		величины. Представить таблицу для группированного статистического ряда
		выборки.
		\item На том же графике представить полигон приведенных частот.
		\item Подобрать параметры функции плотности вероятности. Построить график функции
		плотности вероятности.
		\item Используя критерий согласия Пирсона, проверить гипотезу о принадлежности
		выборки нормально распределенной генеральной совокупности на уровне
		значимости $\alpha = 0.01$. Привести таблицу для критерия $\chi^2$.
		\item Привести график экспериментальной и теоретической функций распределения.
	\end{enumerate}
	
	\quad \textbf{Решение.} Для начала определим значения оценок математического ожидания и дисперсии для данной выборки:
	
	\begin{center}
		$ \hat{\mu}= \dfrac{1}{n} \sum\limits_{i=1}^n {x_i} = -14.344491 $
		
		$ x_{max} = -11.5851, \text{" "}x_{min} = -17.0945 $
		
		$ \hat{\sigma^2} = \dfrac{1}{n-1} \sum\limits_{i=1}^{n} (x_i - \left< x \right>)^2 = 1.21749222668586 $
		
		\begin{table}[h!]
			\centering
			\begin{tabular}{|c|c|c|c|c|c|c|}
				\hline
				$x$ & $(\langle x \rangle -x_i)$ & $(\langle x \rangle-x_i)^2$ \\
				\hline
				-14.1839 & 0.160591 & 0.025789 \\
				-14.2682 & 0.076291 & 0.005820 \\
				$\vdots$ & $\vdots$ & $\vdots$ \\
				-14.4503 & -0.105809 & 0.011196 \\
				-14.4349 & -0.090409 & 0.008174 \\
				-16.0387 & -1.694209 & 2.870344 \\
				\hline
				$\frac{1}{n} \sum\limits_{i=1}^n {x_i} = -14.344491$ & $\sum\limits_{i=1}^n (x_i-\left< x \right>) \approx 0$ & $\frac{1}{n-1} \sum\limits_{i=1}^{n} (x_i - \left< x \right>)^2 = 1.21749222668586$ \\
				\hline
			\end{tabular}
			\caption{определение математического ожидания, линейного отклонения, дисперсии}
		\end{table}
		
	\end{center}
	
	\quad После опредления основных параметров генеральной совокупности необходимо построить гисторгамму относительных частот. Для этого определим оптимальную ширину канала и посчитаем частоту попадания в канал. Согласно лекционному материалу оптимальное количество каналов примерно равно $\sqrt{n}$, тогда для данной выборки количество каналов гистограммы будет равно 10. Ширина канала определяется как: 
	
	\begin{center}
		$ h = \dfrac{x_{max}-x_{min}}{10} = \dfrac{-11.5851 + 17.0945}{10} =  0.55094 $
	\end{center}
	
	\begin{table}[h!]
		\centering
		\begin{tabular}{|c|c|c|c|c|}
			\hline
			№ & [$x_{k-1}, x_k$) & $n_i$ & $ \frac{n_i}{n}$ & $\frac{n_i}{nh}$  \\
			\hline
			1 & [-17.0945, -16.5436) & 3 & 0.03  & 	0.05445783569898719 \\
			
			2 & [-16.5436, -15.9926) & 3 &  0.03 & 0.05445783569898719 \\
			
			3 & [-15.9926, -15.4417) & 9 & 0.09 & 0.16337350709696155 \\
			
			4 & [-15.4417, -14.8908) & 13 & 0.13 & 0.23596035865974516, \\
			
			5 & [-14.8908, -14.3398) & 23 & 0.23 &0.41746832685954915 \\
			
			6 & [-14.3398, -13.7889) & 22& 0.22 &0.3993175300395687 \\
			
			7 & [-13.7889, -13.2380) & 10 & 0.1 &0.18150796819980397, \\
			
			8 & [-13.2380, -12.6871) & 9 & 0.09 &0.16335717137982356 \\
			
			9 & [-12.6871, -12.1361) & 6& 0.06&0.10890478091988239 \\
			
			10 & [-12.1361, -11.5851] & 2 & 0.02 & 0.036301594 \\
			\hline
		\end{tabular}
		\caption{построение гистрограммы относительных частот}
	\end{table}
	
	\begin{figure}[h!]
		\centering
		\includegraphics[width=1\textwidth]{../images/гистограмма_2.png}
		\caption{ Гистограмма относительных частот, полигон относительных частот, функция плотности вероятности.}
	\end{figure}
	
	
	\quad Теперь построим функцию плотности вероятности для нормального распределения. Функция плотности веротяности имеет вид:
	
	\begin{center}
		$ f(x) = \dfrac{1}{\sqrt{2\pi}\sigma} e^{-\frac{(x - \mu)^2}{2\sigma^2}}$
	\end{center}

	\quad Проверим гипотезу о нормальном распределении генеральной совокупности. Для этого воспользуемся критерием $\chi^2$. Теоритические частоты определяются как: 
	
	\begin{center}
		$ n_i ' =  \displaystyle\int\limits_{x_i}^{x_{i+1}} {\dfrac{1}{\sqrt{2\pi}\sigma} e^{\frac{-(x-\mu)^2}{\sigma^2}}} dx $
	\end{center}
	
	\quad Для нахождения критического значения $\chi^2_{\alpha}$ необходимо определить число степеней свободы. В данном случае оно  определяется как количество каналов гистограммы относительных часот за вычетом количества параметров в распределении:
	
	\begin{center}
		
		$ k = \sqrt{n} - 3 = 10 - 3 = 7 $
		
		$ \chi^2 = \sum\limits_{i=1}^{n} \dfrac{(n_i ' - n_i)^2}{n_i '} \approx 4.64$
	\end{center}
	
	\quad $ \implies $ Теперь сравним полученное значение с критическим: $\chi^2_\alpha = 18.5 > 4.64 \implies $ Гипотеза принимается на данном уровне значимости.
	
	\begin{figure}[h!]
		\centering
		\includegraphics[width=1\textwidth]{../images/rasp_2.png}
		\caption{ Теоритическая и эксперементальная функции распределения вероятности.}
	\end{figure}
	
	\begin{center}
		\begin{tabular}{|c|c|c|c|c|c|}
			\hline
			$\frac{(x_i - \mu)}{\sigma}$ & $\frac{(x_{i+1} - \mu)}{\sigma}$ & $\frac{n'_i}{n}$ & $n'_i$ & $n_i$ & $\chi_\alpha^2$ \\
			\hline
			$-\inf$ & -1.9930 & 0.0231 & 2.31 & 3 & 0.20610389610389607\\
			\hline
			-1.9930 & -1.4937 & 0.0445 & 4.45 & 3 & 0.4724719101123596\\
			\hline
			-1.4937 & -0.9944 & 0.0924 & 9.24 & 9 &0.006233766233766245 \\
			\hline
			-0.9944 & -0.4950 & 0.1502 & 15.02 & 13 & 0.27166444740346196 \\
			\hline
			-0.4950 & 0.0043 & 0.1914 & 19.14 & 23 & 0.7784535005224658 \\
			\hline
			0.0043 & 0.5036 & 0.1910 & 19.10 & 22 & 0.44031413612565395\\
			\hline
			0.5036 & 1.0029 & 0.1493 & 14.93 & 10 & 1.627923643670462 \\
			\hline
			1.0029 & 1.5022 & 0.0914 & 9.14 & 9 & 0.0021444201312910455 \\
			\hline
			1.5022 & 2.0015 & 0.0439 & 4.39 & 6 & 0.5904555808656039\\
			\hline
			2.0015 & $+\inf$ & 0.0226 & 2.26 & 2  & 0.02991150442477871\\
			\hline
			\multicolumn{3}{|r|}{} & $\Sigma = 99.98$ & $\Sigma = 100$ & $\Sigma = 4.638066186124713$ \\
			\hline
		
		\end{tabular}
	\end{center}
	
	
	
	\quad \textbf{Выводы:} проведены исселедования двух генеральных совокупностей, найдены их главные, характеризующие значения (дисперсия, мат. ожидание). Проверены гипотезы о распределении этих генеральных совокупностей с помощью критерия Пирсона. В итоге первая выборка распределена по равномерному закону, а вторая по нормальному. 

	
	
	
\end{document}